P14.3 EXACT LIVE PROOF PACKET

BEGIN EXACT LIVE CHUNK
source: manuscript/parts/part5/ch60_mp_psh.tex
label: prop:mp-psh-decreasing-stability
chunk-sha256: a716ec32bb57c705c97dddba3ba047b7a08c7ee75bf8659763406809256c787c
BEGIN CHUNK BYTES
\begin{proof}
Put \(K=X^{\bethan}\) and
\(\mathcal H=\PL(K)\cap\PSH(\alpha)\).  For every \(j\), choose a
decreasing net
\[
                    h_{j,\lambda}\downarrow\phi_j,
                    \qquad h_{j,\lambda}\in\mathcal H.
\]
Let \(\mathcal S\) be the set of all \(h\in\mathcal H\) for which
\[
                  h\geq\phi_j+\epsilon
\tag{60.16a}\label{eq:mp-strict-model-dominant}
\]
for some \(j\geq1\) and some \(\epsilon\in\mathbf Q_{>0}\).  This set is
nonempty because rational translation preserves the model psh cone.

We claim that \(\mathcal S\) is downward directed for pointwise order.
Suppose that \(h_i\geq\phi_{j_i}+\epsilon_i\), \(i=1,2\), as in
\eqref{eq:mp-strict-model-dominant}.  Choose
\(j\geq\max\{j_1,j_2\}\) and a positive rational number
\[
                  4\delta<\min\{\epsilon_1,\epsilon_2\}.
\]
The continuous function \(g=\min\{h_1,h_2\}\) then satisfies
\(\phi_j+4\delta<g\) on the compact space \(K\).  The increasing open
sets
\[
                  \{h_{j,\lambda}+2\delta<g\}
\]
cover \(K\).  Compactness, followed by directedness of the approximating
net, gives one index \(\lambda\) for which this set is all of \(K\).
Hence
\[
                  h:=h_{j,\lambda}+\delta\in\mathcal S,
                  \qquad h\leq h_1,\ h_2.
\]
The minimum \(g\) is only a continuous comparison function; no closure of
the psh cone under minima is being used.

Order \(\mathcal S\) by reverse pointwise order.  The identity net is then
a decreasing net in \(\mathcal H\).  Every member of \(\mathcal S\)
dominates \(\phi\), so its infimum is at least \(\phi\).  Conversely, fix
\(x\in K\) and \(b>\phi(x)\).  Choose \(j\) with \(\phi_j(x)<b\), then
choose \(\epsilon\in\mathbf Q_{>0}\) and \(\lambda\) so that
\(h_{j,\lambda}(x)+\epsilon<b\).  Then
\(h_{j,\lambda}+\epsilon\in\mathcal S\) is smaller than \(b\) at
\(x\).  This proves
\[
                         \inf_{h\in\mathcal S}h=\phi
\]
pointwise, including at points where \(\phi=-\infty\).  Since the limit is
not identically \(-\infty\), \Cref{def:mp-general-psh} gives
\(\phi\in\PSH(\alpha)\).
\end{proof}
END CHUNK BYTES
END EXACT LIVE CHUNK

BEGIN EXACT LIVE CHUNK
source: manuscript/parts/part6/ch71_pseudoeffective_mp_dxz.tex
label: thm:p143-mp-section
chunk-sha256: a800b4975c91d74da5d55bc7bcfdbc39b7d164628da6421e7a6cbeaf1182e109
BEGIN CHUNK BYTES
\begin{proof}
Fix \(\phi\in\PSH_{\mathrm{MP}}(\alpha)\).  By
\Cref{def:mp-general-psh}, there is a decreasing net
\[
 \phi_i\in\PL(X^{\bethan})\cap\PSH_{\mathrm{MP}}(\alpha),
 \qquad \phi_i\downarrow\phi.
\]
Every rational PL MP-psh function has finite model energy, so
\eqref{eq:pseff-finite-energy-bijection} supplies a unique
\(\boldsymbol\Gamma^i\in
\mathcal E^{1,\NA}_{\mathrm{DXZ}}(X,\omega)\).  By
\Cref{cor:p143-all-point-kahler-comparison},
\[
                 \mathfrak T_{\boldsymbol\Gamma^i}=\phi_i
                 \quad\text{on }X^{\bethan}.
\]
If \(i\geq j\), then \(\phi_i\leq\phi_j\) at the positive rational
divisorial points and at the trivial semivaluation.  The order-reflection
statement \eqref{eq:pseff-order-reflection} therefore gives
\(\boldsymbol\Gamma^i\leq\boldsymbol\Gamma^j\).  Thus the DXZ net is
decreasing; no order-compatibility hypothesis has been added.

The trivial semivaluation belongs to \(X^{\divis}\), and
\Cref{thm:mp-psh-divisorial} gives \(\phi(0)>-\infty\).  The endpoint
identity in \Cref{thm:pseff-perturbation-independent-transform} yields
\[
 \boldsymbol\Gamma^i_{\max}
   =\mathfrak T_{\boldsymbol\Gamma^i}(0)
   =\phi_i(0)\downarrow\phi(0)>-\infty.
\]
Hence the decreasing infimum
\(\boldsymbol\Gamma:=\inf_i\boldsymbol\Gamma^i\) exists in
\(\PSH^{\NA}_{\mathrm{DXZ}}(X,\omega)\).  The standard endpoint extension
in \Cref{def:operations-infimum} gives
\[
 \boldsymbol\Gamma_{\max}
   =\inf_i\boldsymbol\Gamma^i_{\max}
   =\phi(0).
\]
Therefore \Cref{thm:pseff-perturbation-independent-transform} gives
\(\mathfrak T_{\boldsymbol\Gamma}(0)=\phi(0)\).  At every positive
rational divisorial point \(q\), \Cref{prop:pseff-operation-formulas} gives
\[
 \mathfrak T_{\boldsymbol\Gamma}(q)
   =\inf_i\mathfrak T_{\boldsymbol\Gamma^i}(q)
   =\inf_i\phi_i(q)
   =\phi(q).
\]
Thus the two functions agree on all of \(X^{\divis}\).

By
\Cref{thm:p142-raw-transform-usc,cor:p142-pseudoeffective-raw-usc},
\(\mathfrak T_{\boldsymbol\Gamma}\) is upper semicontinuous.  It is
antitone for the ideal order: if \(q\preceq v\), then every multiplier-ideal
evaluation at \(q\) is at most the corresponding evaluation at \(v\), and
the defining supremum for the raw transform reverses this inequality.
The comparison clause of \Cref{thm:p143-dominated-divisorial-density},
applied to \(F=\mathfrak T_{\boldsymbol\Gamma}\) and the MP-psh function
\(\phi\), now promotes the divisorial equality to
\[
                 \mathfrak T_{\boldsymbol\Gamma}=\phi
                 \quad\text{on }X^{\bethan}.
\]

Set \(\mathcal S_\alpha(\phi)=\boldsymbol\Gamma\).  If a second DXZ
object has raw transform \(\phi\), the injectivity in
\Cref{thm:pseff-divisorial-inversion} identifies it with
\(\boldsymbol\Gamma\); this proves both uniqueness and independence of the
approximating model net.  Finally, for \(\phi,\psi\in
\PSH_{\mathrm{MP}}(\alpha)\), the equivalence
\eqref{eq:pseff-order-reflection} applied to their all-point raw transforms
gives
\[
 \mathcal S_\alpha(\phi)\leq\mathcal S_\alpha(\psi)
 \quad\Longleftrightarrow\quad
 \phi\leq\psi.
\]
Thus the section preserves and reflects order.
\end{proof}
END CHUNK BYTES
END EXACT LIVE CHUNK

BEGIN EXACT LIVE CHUNK
source: manuscript/parts/part6/ch71_pseudoeffective_mp_dxz.tex
label: thm:p143-full-kahler-comparison
chunk-sha256: 099f110a3c0e429f8efa3d4b7580ab4ba1aa67eb2bf8aa82524337972e1a6ac8
BEGIN CHUNK BYTES
\begin{proof}
Fix \(\boldsymbol\Gamma\), write
\(c=\boldsymbol\Gamma_{\max}\), and set
\(c_j=\max\{c,-j\}\).  The constant object in
\eqref{eq:p143-dxz-truncation} is the one constructed in
\Cref{eq:pseff-constant-transform}.  In every K\"ahler perturbation class
\(\omega+\eta\), its level is zero for \(\tau<-j\).  The maximum
\(\boldsymbol\Gamma^{(j)}\) dominates this constant object.  Since its
\(\mathcal I\)-model levels are normalized to be at most zero, one has
\[
       \bigl(\boldsymbol\Gamma^{(j),\omega+\eta}\bigr)_\tau=0
       \qquad(\tau<-j).
\tag{71.41m}\label{eq:p143-truncation-zero-tail}
\]
Consequently its \(-\infty\)-level has full mass in every perturbed
component and
\[
             \vol(\boldsymbol\Gamma^{(j)})
             =V_\alpha:=\int_X\omega^n>0.
\]
The positive-volume recovery theorem
\Cref{thm:tna-positive-volume} therefore represents
\(\boldsymbol\Gamma^{(j)}\) by an unperturbed positive-mass test curve
\(\gamma^{(j)}\) in the class \(\omega\).  The recovery formula
\eqref{eq:tna-positive-recovery} and
\eqref{eq:p143-truncation-zero-tail} give
\(\gamma^{(j)}_\tau=0\) for \(\tau<-j\).  Hence
\[
 \mathbf E^0(\gamma^{(j)})
 =c_jV_\alpha+
   \int_{-j}^{c_j}
      \bigl(M_\omega(\gamma^{(j)}_\tau)-V_\alpha\bigr)\,\rd\tau
 \in[-jV_\alpha,c_jV_\alpha].
\tag{71.41n}\label{eq:p143-truncation-energy-bound}
\]
Thus \(\boldsymbol\Gamma^{(j)}\) has finite energy.  Notice that this
argument uses the actual zero tail of the recovered curve; it does not
infer finite energy from a merely divisorial lower bound.

Maximum is monotone, so the objects in
\eqref{eq:p143-dxz-truncation} decrease with \(j\).  Their endpoints are
\(c_j\downarrow c\), hence their decreasing infimum exists.  At every
positive rational divisorial point \(q\),
\Cref{prop:pseff-operation-formulas} gives
\[
 \begin{split}
 \mathfrak T_{\inf_j\boldsymbol\Gamma^{(j)}}(q)
   &=\inf_j\max\{\mathfrak T_{\boldsymbol\Gamma}(q),-j\}\\
   &=\mathfrak T_{\boldsymbol\Gamma}(q).
 \end{split}
\]
The endpoint formula gives the same equality at the trivial
semivaluation.  Injectivity in
\Cref{thm:pseff-divisorial-inversion} therefore yields
\[
                         \inf_j\boldsymbol\Gamma^{(j)}
                         =\boldsymbol\Gamma.
\tag{71.41o}\label{eq:p143-dxz-truncation-limit}
\]

Let \(\phi_j\in\mathcal E^1_{\mathrm{MP}}(\alpha)\) correspond to
\(\boldsymbol\Gamma^{(j)}\) under
\eqref{eq:pseff-finite-energy-bijection}.  By
\Cref{cor:p143-all-point-kahler-comparison},
\[
                    \phi_j
                    =\mathfrak T_{\boldsymbol\Gamma^{(j)}}
                    \quad\text{on }X^{\bethan}.
\]
The sequence \((\phi_j)\) is decreasing by monotonicity of the raw
transform.  Moreover,
\[
                    \inf_j\phi_j(0)=\inf_jc_j=c> -\infty,
\]
so \(\phi:=\inf_j\phi_j\) is not identically \(-\infty\).
\Cref{prop:mp-psh-decreasing-stability} gives
\(\phi\in\PSH_{\mathrm{MP}}(\alpha)\).  The divisorial maximum formula
and the endpoint formula show
\[
                    \phi=\mathfrak T_{\boldsymbol\Gamma}
                    \quad\text{on }X^{\divis}.
\]
By
\Cref{thm:p142-raw-transform-usc,cor:p142-pseudoeffective-raw-usc},
\(\mathfrak T_{\boldsymbol\Gamma}\) is upper semicontinuous.  It is
antitone for the ideal order by its defining multiplier-ideal formula.
The comparison clause of
\Cref{thm:p143-dominated-divisorial-density} now promotes the preceding
equality to all of \(X^{\bethan}\).  This proves
\eqref{eq:p143-raw-truncation-limit} and shows that the raw transform maps
the entire DXZ cone into the MP cone.

\Cref{thm:p143-mp-section} already gives
\(\mathfrak T_\alpha\circ\mathcal S_\alpha=\mathrm{id}\).  The raw
transform is injective by
\Cref{thm:pseff-divisorial-inversion}; hence
\(\mathcal S_\alpha\circ\mathfrak T_\alpha=\mathrm{id}\) as well.
Pointwise monotonicity of the raw transform and
\eqref{eq:pseff-order-reflection} show that both maps preserve and reflect
order.  Finally, on the finite-energy locus the divisorial identity
\eqref{eq:pseff-finite-energy-divisorial-link} and injectivity identify
this map with \Cref{thm:pseff-kahler-finite-energy-link}; its energy
identity therefore applies.
\end{proof}
END CHUNK BYTES
END EXACT LIVE CHUNK
